\section{Conclusion}
\label{sec:conclusion}

Typed decision models return a verdict with a probability for every input and output, cheaply enough to make universal screening affordable.
We have derived seven testable requirements from the \pol{} specification and assessed the first weeks of evidence on these models against them.
PoL governance under this text asks more of typed models than the evidence supports.
Hosted \jev{} and the typed models tested beside it should not adjudicate disputes, verify ethical alignment automatically, quantify persons or make a block of a person's speech final.
For quantifying persons, the verdict also rests on precaution, since no study measures error by group.
As a precaution, not as a finding, the verdict extends to the low-cost generative evaluators tested beside them, which are not shown to escape these failures (most were never compared).

For the functions that detect (the safety valve, output truncation and anomaly detection on text), hosted \jev{} is usable only as a recorded, three-valued, locally calibrated detector.
The design in which it is usable requires agreement between independently built detectors before any automatic action, keeps any restriction of a person brief and reversible, ends every chain that could block with a person, audits confident decisions and records every decision.
The public-decision assessment of \clause{5.6} is usable only as an audited aggregate measurement.
Every finding behind this verdict has low or very low certainty, four are contested, and the verdict extrapolates from English benchmarks to a three-valued construct in Chinese that no study has measured (\cref{sec:limitations}).

Several of \pol{}'s own commitments, listed in \cref{sec:verdict}, already match what the evidence recommends.
The evidence turns other provisions into choices the text has yet to make.
The text has to decide whether to draw an explicit line between understanding emotion in language and inferring persons' emotional states, how to ask three-valued questions of models that say ``insufficient'' unreliably, and how to reconcile transparency with closed models.
It must also decide whether non-intervention should yield, for consequential decisions, to a human who can suspend them, and whether per-person assessment can be squared with equal connection.
The next step is to make these choices explicit and to test the resulting safeguard on the specification's own construct.
